\documentclass[12pt]{article}
\usepackage[utf8]{inputenc}
\usepackage[T1]{fontenc}
\usepackage{hyperref}
\title{Alu Kurumba Language Resources for K-12 Curriculum Development}
\author{Compiled by Claude Sonnet 5.5 (claude-sonnet-5-5) for Andr\'as Kornai}
\date{}
\begin{document}
\maketitle

\section{Introduction}
Alu Kurumba (ISO 639-3 xua) is a Dravidian language of the Nilgiri region of India. It is sometimes treated as a dialect of Kannada. Speaker estimates in the search results vary: about 2,400 in ASJP and about 1,000 in the Endangered Languages Project. This survey is based on one web search, and I did not open the individual pages.

\section{Primary Educational Resources}
No primers, textbooks or school materials were found.

\section{Linguistic and Reference Materials}
\begin{itemize}
\item Kapp (1982), \emph{Die Sprache der Alu-Kurumbas} (Harrassowitz, Wiesbaden), which the search summary describes as containing a grammar, texts and a dictionary. It is cited on the ASJP page. \url{https://asjp.clld.org/languages/KURUMBA_ALU}
\item ASJP wordlist for Alu Kurumba (compiled by Agustina Carrizo). \url{https://asjp.clld.org/languages/KURUMBA_ALU}
\item Wikipedia, Alu Kurumba language. \url{https://en.wikipedia.org/wiki/Alu_Kurumba_language}
\item Wikipedia, Kurumba languages. \url{https://en.wikipedia.org/wiki/Kurumba_languages}
\item Endangered Languages Project, Alu Kurumba (bibliography). \url{https://www.endangeredlanguages.com/lang/3567/bibliography}
\end{itemize}

\section{Religious and Community Resources}
\begin{itemize}
\item Scripture Earth, Alu Kurumba (Gospel of Mark and audio recordings, per the search summary). \url{https://www.scriptureearth.org/00eng.php?iso=xua}
\end{itemize}

\section{Digital and Audio Resources}
\begin{itemize}
\item DoBES (Max Planck Institute for Psycholinguistics), Kurumba language documentation project, which covers Alu Kurumba among four groups. \url{https://dobes.mpi.nl/projects/kurumba/project/?lang=id}
\end{itemize}

\section{Curriculum Development Context}
The language is endangered and is not known to be taught in schools. The DoBES documentation and Kapp's 1982 work are the most substantial resources found. Any curriculum would need to be developed with the community.
\end{document}